% CV of Amartya Yadav
%
% Build:  npm run cv        (or: cd cv && latexmk -pdf amartya-yadav-cv.tex)
% Output: cv/amartya-yadav-cv.pdf, linked from the website.
%
% Lines marked "TODO" are placeholders to fill in as things happen.

\documentclass[10pt,a4paper]{article}

\usepackage[T1]{fontenc}
\usepackage[utf8]{inputenc}
\usepackage{lmodern}
\usepackage[a4paper,margin=2cm,top=1.6cm,bottom=1.8cm]{geometry}
\usepackage{xcolor}
\usepackage{titlesec}
\usepackage{enumitem}
\usepackage{tabularx}
\usepackage[hidelinks]{hyperref}

\definecolor{accent}{RGB}{153,0,0}

\pagestyle{empty}
\setlength{\parindent}{0pt}
\setlist[itemize]{leftmargin=1.2em,itemsep=1pt,topsep=2pt,parsep=0pt}

\titleformat{\section}{\large\bfseries\color{accent}}{}{0pt}{}[\vspace{-6pt}\rule{\linewidth}{0.4pt}]
\titlespacing*{\section}{0pt}{10pt}{6pt}

% \entry{title}{dates}{organisation}{location}
\newcommand{\entry}[4]{%
  \textbf{#1}\hfill #2\\
  \textit{#3}\hfill \textit{#4}\par\vspace{2pt}}

\hypersetup{
  pdftitle={Amartya Yadav -- CV},
  pdfauthor={Amartya Yadav}
}

\begin{document}

% ------------------------------------------------------------------ header
{\centering
  {\LARGE\bfseries Amartya Yadav}\\[4pt]
  PhD Research Fellow in High-Performance Computing (from 1 December 2026)\\
  University of Oslo \& Simula Research Laboratory\\[4pt]
  \href{mailto:amartyadav@gmail.com}{amartyadav@gmail.com} \;\textbar\;
  \href{https://amartyadav.com}{amartyadav.com} \;\textbar\;
  \href{https://codeberg.org/amartya}{codeberg.org/amartya} \;\textbar\;
  \href{https://github.com/amartyadav}{github.com/amartyadav}\par
}

% ------------------------------------------------------------------ interests
\section*{Research Interests}
Convergence of HPC and AI; numerical methods and parallel solvers (algebraic multigrid, Gauss--Seidel);
computational fluid dynamics (Lattice Boltzmann Method); wafer-scale and unconventional architectures;
number formats (IEEE~754 vs.\ posits); performance engineering.

% ------------------------------------------------------------------ positions
\section*{Positions}

\entry{PhD Research Fellow, High-Performance Computing}{From Dec 2026}
      {Dept.\ of Informatics, University of Oslo \& Simula Research Laboratory}{Oslo, Norway}
\begin{itemize}
  \item Full-time, three-year research position with Prof.\ Xing Cai on the convergence of HPC and AI.
  % TODO: add project details once the PhD is under way.
\end{itemize}

\entry{HPC Software Engineer}{Feb 2026 -- Sep 2026}
      {Calligo Technologies}{Bengaluru, India}
\begin{itemize}
  \item Implemented, profiled and compared parallel solvers on x86 (IEEE~754) and the UTTUNGA RISC-V
        board with experimental posit arithmetic.
  \item Multi-material thermal analysis (Poisson, HYPRE BoomerAMG) of a CFRP + aluminium domain. The posit
        version on UTTUNGA converged in about 11--12 iterations at tolerances of $10^{-14}$--$10^{-15}$, where x86 failed to converge.
  \item Added OpenMP red-black Gauss--Seidel smoothing to OpenFOAM. Profiled with AMD uProf, Intel VTune,
        Intel Advisor and LIKWID (cache, memory bandwidth, roofline, power).
  \item Built a Lattice Boltzmann CFD kernel with a physical validation suite, in serial and OpenMP versions at
        single and double precision, benchmarked on UTTUNGA against x86.
  \item Started the company's first FPGA project, porting OpenFOAM hotspot functions to FPGA with HLS.
  \item Wrote a technical report for every project. Published two OpenFOAM plugins on Zenodo.
\end{itemize}

\entry{Freelance Software Developer}{2023 -- 2024}
      {Hospitals in Gorakhpur}{Gorakhpur, India}
% TODO: check these dates
\begin{itemize}
  \item Built electronic patient record and staff management systems.
\end{itemize}

\entry{Software Developer (Full-Stack)}{2022 -- 2023}
      {Fennex}{Aberdeen, UK}
\begin{itemize}
  \item ASP.NET, TypeScript, React, SQL. Did much of the UI design for the flagship BBSS product on iOS and Android.
\end{itemize}

% ------------------------------------------------------------------ education
\section*{Education}

\entry{MSc High Performance Computing, with Merit}{2024 -- 2025}
      {EPCC, University of Edinburgh}{Edinburgh, UK}
\begin{itemize}
  \item Dissertation: Lattice-Boltzmann methods for CFD on the Cerebras CS-3 (Taylor--Green vortex decay in CSL),
        supervised by Dr.\ Justs Zarins.
  \item Courses: Advanced Message Passing Programming; Performance Programming; Accelerated Systems;
        Machine Learning at Scale; Parallel Design Patterns; MPI; OpenMP; Numerical Algorithms in HPC;
        Design and Analysis of Parallel Algorithms; HPC Architectures.
\end{itemize}

\entry{BSc (Hons) Computer Science, 2:1}{2018 -- 2022}
      {Robert Gordon University}{Aberdeen, UK}

% ------------------------------------------------------------------ publications
\section*{Publications, Software \& Talks}

\textbf{Peer-reviewed papers}\\
\textit{None yet. Expected during the PhD.}
% TODO: add papers here, newest first.

\vspace{4pt}
\textbf{Software}
\begin{itemize}
  \item A.~Yadav. \textit{Red-Black Parallel Gauss-Seidel Smoother with DIA Memory Access Pattern}.
        OpenFOAM plugin, Zenodo, April 2026. % TODO: add DOI
  \item A.~Yadav. \textit{Gauss-Seidel Smoother with DIA Memory Access Pattern}.
        OpenFOAM plugin, Zenodo, April 2026. % TODO: add DOI
\end{itemize}

\textbf{Talks \& tutorials}
\begin{itemize}
  \item \textit{Efficient Halo Exchange for Stencil-Based Scientific Codes on Cerebras Wafer-Scale Engine}.
        Tutorial, Supercomputing India 2025, Bengaluru, November 2025.
  \item \textit{Multi-Material Thermal Analysis with HYPRE BoomerAMG and Posits on UTTUNGA}.
        Live demo, investment summit led by Seafund Ventures, June 2026.
\end{itemize}

% ------------------------------------------------------------------ skills
\section*{Technical Skills}
\begin{tabularx}{\linewidth}{@{}l X@{}}
  \textbf{Languages}   & C (primary), C++, Rust, Java, CSL; also Kotlin, Swift, C\#, TypeScript, Python, Elixir \\
  \textbf{Parallel}    & MPI, OpenMP, CUDA, SYCL, Cerebras WSE-3, FPGA (HLS) \\
  \textbf{Numerics}    & Lattice Boltzmann Method, Poisson solvers, HYPRE BoomerAMG, OpenFOAM, posit arithmetic \\
  \textbf{Profiling}   & AMD uProf, Intel VTune, Intel Advisor, LIKWID \\
  \textbf{Visualisation} & ParaView, GLVis \\
  \textbf{Systems}     & Cerebras CS-3, Cirrus (EPCC), UTTUNGA (RISC-V), x86 (AMD/Intel)
\end{tabularx}

% ------------------------------------------------------------------ projects
\section*{Selected Independent Projects}
\begin{itemize}
  \item \textbf{Mandi Mitra} (iOS \& Android, FOSS): grocery price calculator for Indian farmers' markets.
        No ads, no trackers.
  \item \textbf{QuickPanchang} (macOS App Store; Electron): Hindu calendar app, backed by my own scrapers,
        ASP.NET API and caching proxy.
\end{itemize}

% ------------------------------------------------------------------ languages
\section*{Languages}
English, Hindi, Urdu, Bhojpuri, Gujarati, Punjabi (fluent); Norwegian (learning).
Scripts: Devanagari, Latin, Gujarati, Kaithi.

\vfill
{\footnotesize\color{gray} Last updated: \today. Built with \LaTeX; source at amartyadav.com/cv/amartya-yadav-cv.tex}

\end{document}
